\section{Технический проект}

\subsection{Общая характеристика проектного решения}

Система реализована как веб-приложение с клиент-серверной архитектурой \cite{art_client_server}. Пользователь работает с ней из браузера, на компьютер ничего ставить не нужно. На сервере располагаются данные, бизнес-логика и проверки прав. Клиент отвечает за интерфейс и реакцию на действия пользователя.

Клиентская часть -- одностраничное приложение (SPA) на React. Полная загрузка страницы происходит один раз, при первом входе. Дальше браузер обращается к серверу только за конкретными данными, а сам интерфейс не собирается заново при каждом переходе между разделами. Трафика меньше, между разделами нет видимых пауз.

Серверная часть написана на Django с подключенной библиотекой DRF \cite{drf_docs}. Программный интерфейс построен в стиле REST \cite{art_rest, rfc9110}. У каждой основной сущности системы -- проектов и задач, зависимостей между задачами, участников проекта и самих пользователей -- свой набор маршрутов с типовыми операциями: получение, создание, изменение, удаление. Обмен данными идет в формате JSON по HTTP.

Постоянное хранилище -- реляционная СУБД MySQL. Запросы к ней строит ORM-слой Django, никакой сборки SQL из пользовательского ввода в коде нет, и основная поверхность атак через SQL-инъекции при таком подходе закрыта \cite{owasp_top10}.

В системе пять слоев ответственности \cite{fowler_patterns}:
\begin{itemize}
    \item уровень представления -- компоненты React, страницы, маршрутизация, локальное состояние интерфейса;
    \item уровень обмена -- модуль на стороне клиента, который формирует запросы к серверу и разбирает ответы;
    \item уровень программного интерфейса -- представления и сериализаторы DRF, маршруты, фильтры, проверки прав;
    \item уровень бизнес-логики -- модели Django, методы наборов запросов, проверки на стороне приложения, транзакции;
    \item уровень хранения -- таблицы MySQL, индексы, ограничения целостности.
\end{itemize}

Вход в систему реализован на основе базового механизма сеансов в Django. После проверки логина и пароля сервер открывает сеанс и кладет его идентификатор в защищенную cookie. Браузер сам подставляет эту cookie ко всем последующим запросам. Изменяющие операции дополнительно подтверждаются CSRF-токеном, который сервер отдает отдельной cookie.

Общая схема организации системы приведена на рисунке~\ref{fig:arch}.

\begin{figure}[H]
\centering
\includegraphics[height=0.5\textheight]{arch}
\caption{Общая схема организации системы}
\label{fig:arch}
\end{figure}

В локальной разработке сервер поднимается через Uvicorn в ASGI-режиме, клиент -- через сервер разработки React. На промышленном стенде собранный клиентский код раздается как набор статических файлов, запросы к API идут по HTTPS.

\subsection{Выбор средств и технологий реализации}

\subsubsection{Состав применяемого технологического стека}

Для реализации системы выбран следующий набор инструментов:
\begin{itemize}
    \item язык программирования Python версии 3.10 или выше -- серверная часть;
    \item веб-фреймворк Django версии 6.0 -- маршрутизация запросов, модели, сеансы, проверки CSRF;
    \item библиотека Django REST Framework версии 3.16 -- построение программного интерфейса в стиле REST;
    \item библиотека drf-nested-routers -- вложенные маршруты вида /projects/\{id\}/tasks;
    \item библиотека django-filter -- фильтрация коллекций по параметрам запроса;
    \item библиотека django-cors-headers -- обработка кросс-доменных запросов в режиме разработки;
    \item библиотека React версии 19 -- клиентская часть;
    \item библиотека react-router-dom версии 7 -- маршрутизация на стороне клиента;
    \item СУБД MySQL версии 8.0 -- постоянное хранилище;
    \item сервер приложений Uvicorn -- запуск серверного процесса.
\end{itemize}

Стек подбирался с учетом нескольких соображений \cite{orlov_se}. Главным был зрелостный критерий: Django, React и MySQL давно работают на боевых проектах в коммерческих компаниях и к экспериментальным технологиям не относятся, поэтому риск столкнуться с нерешенными проблемами фреймворка минимален. По каждой технологии есть подробная документация и крупное сообщество, что упрощает поиск ответов на вопросы реализации. Лицензионная сторона тоже устроена: все позиции бесплатны и распространяются с открытым исходным кодом. Наконец, между собой эти технологии соединяются через стандартный канал -- HTTP с телами в JSON -- что позволяет заменить любую часть на аналог без переписывания остальных.

\subsubsection{Серверный язык программирования и фреймворк Django}

Серверный язык -- Python \cite{lutz_python}. Под него есть пакеты практически на любую задачу -- от работы с базами данных до сетевого взаимодействия. По скорости Python уступает компилируемым языкам, но в данном случае это не узкое место: основное время уходит на запросы к базе и сериализацию ответов.

Django существует с 2005 года и за это время вобрал в себя почти все, что нужно типичному веб-приложению \cite{django_docs, art_python_web}, в частности:
\begin{itemize}
    \item ORM с миграциями -- модели описываются как классы Python, миграции базы создаются автоматически;
    \item система аутентификации с сеансами и хешированием паролей;
    \item защита от CSRF, XSS и кликджекинга на уровне промежуточных обработчиков;
    \item встроенная проверка входных данных и система форм;
    \item панель администратора для просмотра и редактирования записей.
\end{itemize}

Инфраструктурный код в основном уже написан внутри Django -- разработчику остается прикладная логика.

\subsubsection{Django REST Framework для построения программного интерфейса}

DRF добавляет к Django средства для построения REST-интерфейсов. В разработке пригодились следующие его компоненты:
\begin{itemize}
    \item классы-наборы ModelViewSet -- для типового набора операций над сущностями;
    \item сериализаторы ModelSerializer -- двусторонняя сборка JSON-представлений из моделей Django и обратно;
    \item классы разрешений -- проверка прав на уровне запросов и отдельных действий;
    \item постраничный вывод -- для списков проектов через стандартный класс PageNumberPagination;
    \item фильтрация через DjangoFilterBackend -- набор параметров запроса описывается как класс FilterSet.
\end{itemize}

За счет DRF шаблонного кода становится заметно меньше. Например, операции над проектами реализованы классом ProjectViewSet, унаследованным от ModelViewSet. Из дописанного руками -- переопределенный get\_queryset, который ограничивает выборку доступными пользователю проектами, perform\_create с автоматическим добавлением создателя в участники с ролью владельца, а также perform\_update и perform\_destroy с проверкой роли владельца.

\subsubsection{Клиентская библиотека React}

React \cite{react_docs} -- библиотека, на которой собрана клиентская часть. Среди инструментов для интерфейсов на JavaScript она самая распространенная. Компоненты в коде описаны как функции, состояние и побочные эффекты подключаются через хуки useState, useEffect, useCallback, useMemo. Глобальное состояние, общее для разных страниц -- текущий пользователь, список доступных проектов, цветовая тема -- вынесено в три контекста: AuthContext, ProjectsContext, ThemeContext. Маршрутизация подключена через react-router-dom. Разделы интерфейса привязаны к маршрутам вида /projects/\{id\}/gantt, /users/\{username\}, /settings; переход между ними не вызывает перезагрузку документа.

Стили оформлены как CSS-модули: у каждого компонента свой файл со стилями, классы локальны и не конфликтуют между компонентами. Цвета вынесены в CSS-переменные, поэтому смена темы происходит без пересборки приложения.

\subsubsection{СУБД MySQL}

MySQL -- реляционная СУБД с открытым исходным кодом, одна из самых распространенных среди свободно доступных решений. Из возможностей MySQL 8.0 пригодились проверочные ограничения на уровне таблицы. Они нужны для того, чтобы запретить самозависимость задачи и дублирующие пары зависимостей прямо в схеме базы.

Схема управляется миграциями Django, ручное редактирование таблиц запрещено. Кодировка utf8mb4 поддерживает кириллические заголовки задач и проектов.

\subsection{Архитектурное построение системы}

\subsubsection{Уровневая модель архитектуры}

Код в системе разнесен по уровням, у каждого -- своя зона ответственности. Уровень обращается только к следующему за ним, и за счет этого проще локализовать ошибки и заменять отдельные части. Уровневая схема приведена на рисунке~\ref{fig:levels_arch}.

\begin{figure}[H]
\centering
\includegraphics[height=0.45\textheight]{levels_arch}
\caption{Уровневая схема архитектуры системы}
\label{fig:levels_arch}
\end{figure}

Из пяти уровней первые два работают в браузере, следующие два -- на сервере, нижний представляет собой саму базу данных. Граница между клиентом и сервером проходит по HTTP-вызову. На клиенте уровень обмена упаковывает запрос в JSON и отправляет его на сервер, где уровень программного интерфейса принимает и разбирает этот запрос. Между бизнес-логикой и хранением встроен Django ORM, поэтому ни одна строка прикладного кода не собирает SQL вручную.

\subsubsection{Компоненты клиентской части}

На клиенте код собран в три группы файлов:
\begin{itemize}
    \item компоненты -- многоразовые элементы интерфейса (Button, Input, Modal, Select), а также крупные блоки (GanttChart, Header, Sidebar, TaskCard, ProjectMembers, SettingsPanel);
    \item страницы -- сборочные единицы маршрутов (GanttPage, ProjectPage, ProfilePage, SettingsPage, WelcomePage);
    \item контексты -- хранилища глобального состояния (AuthContext, ProjectsContext, ThemeContext).
\end{itemize}

AuthContext держит данные текущего пользователя и операции входа, выхода и регистрации. Через ProjectsContext страницам доступен список проектов пользователя, идентификатор последнего открытого проекта и операции создания, удаления, обновления. Отдельно стоит ThemeContext, отвечающий за оформление: он переключает режим темы (системный, светлый, темный, пользовательский), хранит набор цветов и применяет их к корневому элементу документа через CSS-переменные.

Большая часть клиентского кода приходится на компонент GanttChart. Он отвечает за всю отрисовку диаграммы -- временной шкалы проекта, полос задач, стрелок зависимостей, линии текущей даты, -- а также обрабатывает перетаскивание. Полосы выкладываются на сетку, шаг которой пересчитывается под выбранный масштаб в диапазоне 50--400~\%. Стрелки зависимостей четырех типов соединяют точки на торцах полос.

\subsubsection{Компоненты серверной части}

Серверное приложение Django поделено на несколько пакетов:
\begin{itemize}
    \item models -- описания таблиц как классов Django: User, Project, Task, TaskDependency, ProjectMember, UserColorPalette;
    \item serializers -- сериализаторы DRF, по одному файлу на сущность;
    \item views -- классы-наборы операций над сущностями;
    \item filters -- классы фильтров для списочных запросов;
    \item validators -- проверки, не привязанные к одной модели (проверка пароля, проверка циклов в графе зависимостей задач);
    \item permissions -- функции и классы проверки прав доступа;
    \item mixins -- общие фрагменты поведения для наборов представлений.
\end{itemize}

Маршруты собираются через DefaultRouter из DRF и NestedDefaultRouter из drf-nested-routers. На верхнем уровне доступны /auth, /users, /projects, /tasks. Под маршрутом проекта вложены /projects/\{id\}/tasks, /projects/\{id\}/dependencies, /projects/\{id\}/members.

Права доступа проверяются на двух уровнях. Глобальный уровень требует аутентификации для всех маршрутов, кроме регистрации, входа и получения CSRF-токена. Уровень проекта смотрит, состоит ли пользователь в проекте и какая у него роль. Логика собрана в функциях is\_member, is\_writer, is\_owner модуля permissions. Каждый набор представлений вызывает их в методе initial -- до того, как выполнится само действие.

\subsubsection{Компоненты уровня доступа к данным}

Все обращения к базе идут через Django ORM \cite{art_orm}. Таблицы описаны классами-наследниками models.Model, поля -- атрибутами класса. Связи задаются через ForeignKey, OneToOneField и обратные отношения с параметром related\_name. Ограничения целостности -- уникальность пар, запрет самозависимости задачи -- прописаны в Meta.constraints и переходят в схему базы как ограничения на уровне таблицы.

Если одна операция меняет несколько таблиц сразу, она оборачивается в transaction.atomic. Так, при создании проекта в одной транзакции выполняются запись самого проекта и добавление автора в таблицу участников с ролью владельца. Сбой на любом шаге откатывает обе записи.

Часто исполняемые запросы дополнены select\_related (для связей по внешнему ключу) и prefetch\_related (для обратных отношений). Эти подсказки ORM меняют стратегию выборки: вместо отдельного запроса за участниками на каждый проект из списка и за зависимостями на каждую задачу выполняется одно объединение, покрывающее всю коллекцию.

\subsubsection{Диаграмма последовательности обработки запроса}

Последовательность обработки пользовательского действия показана на рисунке~\ref{fig:tp_seq}. Диаграмма последовательности построена в нотации UML \cite{fowler_uml}. На ней отдельно выделены браузер, React-приложение, серверный слой Django + DRF, ORM и СУБД MySQL -- между этими участниками меняется форма представления данных. Событие интерфейса превращается в HTTP-запрос, затем в операции над моделями и SQL-запросы, а после выполнения возвращается клиенту как JSON-ответ.

\begin{figure}[H]
\centering
\includegraphics[width=\textwidth]{seq_request}
\caption{Последовательность обработки запроса пользователя}
\label{fig:tp_seq}
\end{figure}

Шаги диаграммы:
\begin{enumerate}
    \item Пользователь выполняет действие в браузере: открывает список, сохраняет карточку задачи, меняет порядок работ или создает зависимость.
    \item Браузер передает событие React-приложению. Клиентский код обрабатывает событие, собирает параметры операции и формирует HTTP-запрос с телом в формате JSON.
    \item Запрос отправляется в серверную часть Django + DRF. Для изменяющих операций браузер прикладывает cookie сеанса и CSRF-токен, полученные при входе в систему.
    \item Сервер проверяет сеанс пользователя и корректность CSRF-токена. Если сеанс отсутствует или токен неверен, дальнейшая обработка прекращается.
    \item После аутентификации выполняется проверка прав на проект: сервер определяет, входит ли пользователь в состав проекта и достаточно ли его роли для выбранного действия.
    \item DRF-сериализатор и прикладные валидаторы проверяют входные данные. На этом этапе отбрасываются некорректные типы, отсутствующие обязательные поля, недопустимые роли и нарушения правил предметной области.
    \item Если запрос требует чтения или изменения данных, сервер обращается к ORM. ORM строит SQL-запросы по моделям Django и передает их в MySQL.
    \item MySQL выполняет SQL-запрос и возвращает набор записей. ORM преобразует результат в объекты моделей или QuerySet, с которыми работает серверный код.
    \item DRF сериализует полученные объекты в JSON и возвращает HTTP-ответ React-приложению.
    \item React-приложение обновляет локальное состояние интерфейса. Браузер перерисовывает изменившиеся элементы страницы и показывает пользователю результат операции без полной перезагрузки документа.
\end{enumerate}

\subsubsection{Форматы обмена данными}

Тела запросов и ответов -- всегда JSON. Для календарных полей клиент передает дату строкой в формате <<ГГГГ-ММ-ДД>>. Сервер хранит эти значения как DateTimeField и в ответах возвращает ISO-строки -- клиент использует из них только календарную часть. Идентификаторы записей передаются целыми числами.

Списочные ресурсы приходят страницами по 25 записей. Формат страницы стандартный для DRF -- объект с полями count, next, previous, results, где results -- массив записей текущей страницы, а next и previous -- ссылки на соседние.

При ошибках сервер возвращает JSON-объект с полем detail или словарь с пояснениями по полям. HTTP-коды соответствуют семантике REST: 200 -- успешное чтение, 201 -- создание, 204 -- успешное действие без тела ответа, 400 -- ошибка валидации, 401 -- нет аутентификации, 403 -- недостаточно прав, 404 -- ресурс не найден.

\subsubsection{Аутентификация и авторизация}

Аутентификация работает через стандартный механизм сеансов Django. Сценарий входа выглядит так:
\begin{enumerate}
    \item Перед вызовом изменяющих операций клиент запрашивает CSRF-токен по маршруту /auth/csrf. Сервер ставит cookie csrftoken.
    \item Клиент отправляет логин и пароль на /auth/login. Сервер вызывает функцию authenticate Django, проверяет хеш пароля, открывает сеанс работы и устанавливает cookie sessionid.
    \item При каждом последующем запросе браузер прикладывает обе cookie. Для изменяющих методов клиент дополнительно передает значение CSRF-токена в заголовке X-CSRFToken.
\end{enumerate}

Пароли хранятся в виде хеша. По умолчанию Django вычисляет их через PBKDF2 с солью \cite{art_hashing}. Открытый пароль нигде не сохраняется.

Проверка прав на ресурс делится на два этапа. Сначала DRF убеждается, что у запроса есть действующий сеанс, иначе возвращает 401. После этого серверный код проверяет, состоит ли пользователь в нужном проекте и какая у него роль. От роли зависит набор разрешенных действий: владельцу доступны все операции, включая редактирование самого проекта, удаление его и управление составом участников. У редактора есть права на задачи и зависимости, но менять состав команды он не может. Наблюдатель в проект только смотрит.

Сложность пароля контролирует собственный валидатор PasswordStrengthValidator. Минимальная длина -- восемь символов. В пароле должны быть хотя бы одна строчная латинская буква, одна заглавная, одна цифра и один спецсимвол из заранее заданного набора. Кириллица и пробелы запрещены.

\subsubsection{Описание программного интерфейса}

Программный интерфейс размещен под префиксом /api и не использует завершающий слеш в маршрутах. Ниже приведены основные эндпоинты, тела запросов и ответы сервера. Для изменяющих методов POST, PATCH, PUT и DELETE браузер передает cookie сеанса и заголовок X-CSRFToken. В примерах показаны значимые поля JSON. Списочные запросы по ресурсам users, projects и вложенному списку tasks проекта возвращаются в стандартной постраничной обертке DRF с полями count, next, previous и results по 25 записей на странице. Списки зависимостей задач, участников проекта и задач пользователя (GET /api/tasks) приходят целиком в виде массива без пагинации -- объем этих коллекций ограничен размером отдельного проекта.

\textbf{Получение CSRF-токена.}

GET /api/auth/csrf

Запрос не содержит тела. Сервер устанавливает cookie csrftoken и возвращает код 204 No Content.

\textbf{Вход пользователя.}

POST /api/auth/login

Тело запроса:

\begin{lstlisting}[numbers=none, language=json]
{
  "username": "ivanov.i",
  "password": "Password1!"
}
\end{lstlisting}


Ответ 200 OK:

\begin{lstlisting}[numbers=none, language=json]
{
  "id": 3,
  "full_name": "Иванов Иван Иванович",
  "first_name": "Иван",
  "middle_name": "Иванович",
  "last_name": "Иванов",
  "username": "ivanov.i",
  "email": "ivanov.i@example.com",
  "color_palette": {
    "theme_mode": "system",
    "colors": {
      "appBackground": "#fbf4ff",
      "surface": "#fbf4ff",
      "headerBackground": "#9596af",
      "activeMenuBackground": "#9596af",
      "border": "#9596af",
      "primaryText": "#61638B",
      "secondaryText": "#9596af",
      "activeMenuText": "#fbf4ff",
      "inverseText": "#fbf4ff",
      "accent": "#4193ff",
      "task": "#4193ff",
      "ganttGridLine": "#d8d8e6",
      "ganttCurrentTime": "#f04438",
      "ganttArrow": "#9596af",
      "error": "#e54e4e",
      "danger": "#d94040",
      "success": "#4caf50"
    }
  }
}
\end{lstlisting}


\textbf{Выход пользователя.}

POST /api/auth/logout

Запрос не содержит тела. Сервер завершает сеанс и возвращает код 204 No Content.

\textbf{Регистрация пользователя.}

POST /api/users

Создает новую учетную запись пользователя. Тело запроса:

\begin{lstlisting}[numbers=none, language=json]
{
  "first_name": "Иван",
  "last_name": "Иванов",
  "middle_name": "Иванович",
  "username": "ivanov.i",
  "email": "ivanov.i@example.com",
  "password": "Password1!"
}
\end{lstlisting}


Ответ 201 Created совпадает по структуре с ответом входа пользователя: сервер возвращает созданную запись пользователя с вложенной палитрой color\_palette. Поле password в ответ не включается.

\textbf{Получение текущего пользователя.}

GET /api/users/me

Запрос не содержит тела. Ответ 200 OK совпадает по структуре с объектом пользователя из ответа /api/auth/login.

\textbf{Получение и изменение пользователя по id.}

GET /api/users/\{id\}

Запрос не содержит тела. Ответ 200 OK содержит объект пользователя. Обычный пользователь может получить только свою запись, суперпользователь -- любую.

PATCH /api/users/\{id\}

Обновляет поля профиля пользователя. Тело запроса содержит изменяемые поля профиля:

\begin{lstlisting}[numbers=none, language=json]
{
  "first_name": "Иван",
  "last_name": "Петров",
  "middle_name": "Иванович",
  "email": "petrov.i@example.com"
}
\end{lstlisting}


Ответ 200 OK содержит обновленный объект пользователя. Поле username после создания не изменяется.

\textbf{Список и поиск пользователей.}

GET /api/users?search=ив\&page=1

Возвращает постраничный список пользователей. Поддерживается параметр search -- фильтрует по началу полного имени или логина. Запрос не содержит тела. Ответ 200 OK возвращается постранично:

\begin{lstlisting}[numbers=none, language=json]
{
  "count": 1,
  "next": null,
  "previous": null,
  "results": [
    {
      "id": 3,
      "full_name": "Иванов Иван Иванович",
      "first_name": "Иван",
      "middle_name": "Иванович",
      "last_name": "Иванов",
      "username": "ivanov.i",
      "email": "ivanov.i@example.com",
      "color_palette": { "theme_mode": "system", "colors": {} }
    }
  ]
}
\end{lstlisting}


GET /api/users/search?q=иван

Этот маршрут используется при добавлении участников проекта. Запрос не содержит тела. Ответ 200 OK -- массив не более чем из десяти кратких записей:

\begin{lstlisting}[numbers=none, language=json]
[
  {
    "id": 3,
    "username": "ivanov.i",
    "full_name": "Иванов Иван Иванович"
  }
]
\end{lstlisting}


\textbf{Цветовая палитра текущего пользователя.}

GET /api/users/color-palette

Возвращает цветовую палитру текущего пользователя. Запрос не содержит тела. Ответ 200 OK:

\begin{lstlisting}[numbers=none, language=json]
{
  "theme_mode": "custom",
  "colors": {
    "appBackground": "#fbf4ff",
    "surface": "#ffffff",
    "headerBackground": "#9596af",
    "activeMenuBackground": "#9596af",
    "border": "#d8d8e6",
    "primaryText": "#61638B",
    "secondaryText": "#9596af",
    "activeMenuText": "#fbf4ff",
    "inverseText": "#fbf4ff",
    "accent": "#4193ff",
    "task": "#4193ff",
    "ganttGridLine": "#d8d8e6",
    "ganttCurrentTime": "#f04438",
    "ganttArrow": "#9596af",
    "error": "#e54e4e",
    "danger": "#d94040",
    "success": "#4caf50"
  }
}
\end{lstlisting}


PATCH /api/users/color-palette

Обновляет цветовую палитру текущего пользователя. Тело запроса может содержать theme\_mode и любой поднабор цветов:

\begin{lstlisting}[numbers=none, language=json]
{
  "theme_mode": "custom",
  "colors": {
    "appBackground": "#ffffff",
    "task": "#2f80ed",
    "ganttArrow": "#61638B"
  }
}
\end{lstlisting}


Ответ 200 OK содержит всю обновленную палитру.

\textbf{Список проектов.}

GET /api/projects?page=1\&search=разработка\&overdue=false

Запрос не содержит тела. Поддерживаются параметры search, creator\_id, overdue и page. Ответ 200 OK:

\begin{lstlisting}[numbers=none, language=json]
{
  "count": 1,
  "next": null,
  "previous": null,
  "results": [
    {
      "id": 43,
      "title": "Разработка ИС",
      "description": "Описание проекта",
      "start_date": "2026-02-01T00:00:00Z",
      "end_date": "2026-05-15T00:00:00Z",
      "creator_id": 3,
      "current_role": "owner"
    }
  ]
}
\end{lstlisting}


\textbf{Создание проекта.}

POST /api/projects

Создает новый проект. Тело запроса:

\begin{lstlisting}[numbers=none, language=json]
{
  "title": "Разработка ИС",
  "description": "Описание проекта",
  "start_date": "2026-02-01",
  "end_date": "2026-05-15"
}
\end{lstlisting}


Ответ 201 Created содержит объект проекта. Сервер заполняет creator\_id и current\_role, а также создает запись участника с ролью owner для автора проекта.

\textbf{Получение, изменение и удаление проекта.}

GET /api/projects/\{id\}

Запрос не содержит тела. Ответ 200 OK содержит объект проекта.

PATCH /api/projects/\{id\}

Обновляет поля проекта. Тело запроса содержит изменяемые поля проекта:

\begin{lstlisting}[numbers=none, language=json]
{
  "title": "Разработка ИС",
  "end_date": "2026-06-01"
}
\end{lstlisting}


Ответ 200 OK содержит обновленный объект проекта. Метод PUT использует такой же набор полей, но ожидает полную запись проекта. Удаление выполняется запросом DELETE /api/projects/\{id\} без тела и возвращает 204 No Content.

\textbf{Список задач пользователя.}

GET /api/tasks?search=проектирование\&overdue=false

Запрос не содержит тела. Маршрут возвращает задачи, созданные текущим пользователем; для суперпользователя доступен полный список. Ответ 200 OK -- массив задач:

\begin{lstlisting}[numbers=none, language=json]
[
  {
    "id": 7,
    "title": "Проектирование API",
    "description": "Описать REST API",
    "start_date": "2026-02-10T00:00:00Z",
    "end_date": "2026-02-15T00:00:00Z",
    "number": 1,
    "creator_id": 3,
    "project_id": 43
  }
]
\end{lstlisting}


\textbf{Задачи проекта.}

GET /api/projects/\{project\_id\}/tasks?page=1

Запрос не содержит тела. Поддерживаются параметры search, creator\_id, overdue и page. Ответ 200 OK возвращается постранично:

\begin{lstlisting}[numbers=none, language=json]
{
  "count": 1,
  "next": null,
  "previous": null,
  "results": [
    {
      "id": 7,
      "title": "Проектирование API",
      "description": "Описать REST API",
      "start_date": "2026-02-10T00:00:00Z",
      "end_date": "2026-02-15T00:00:00Z",
      "number": 1,
      "creator_id": 3,
      "project_id": 43
    }
  ]
}
\end{lstlisting}


POST /api/projects/\{project\_id\}/tasks

Создает задачу в проекте. Тело запроса:

\begin{lstlisting}[numbers=none, language=json]
{
  "title": "Проектирование API",
  "description": "Описать REST API",
  "start_date": "2026-02-10",
  "end_date": "2026-02-15"
}
\end{lstlisting}


Ответ 201 Created содержит созданную задачу. Поля id, number, creator\_id и project\_id назначаются сервером.

\textbf{Получение, изменение и удаление задачи.}

Задачу можно адресовать по глобальному идентификатору /api/tasks/\{id\} или по номеру внутри проекта /api/projects/\{project\_id\}/tasks/\{number\}. Для GET тело не передается, ответ 200 OK содержит объект задачи.

Для изменения используется PATCH или PUT:

\begin{lstlisting}[numbers=none, language=json]
{
  "title": "Проектирование API",
  "end_date": "2026-02-20"
}
\end{lstlisting}


Ответ 200 OK содержит обновленную задачу. Удаление выполняется запросом DELETE без тела и возвращает 204 No Content.

\textbf{Изменение порядка задач.}

POST /api/projects/\{project\_id\}/tasks/move

Меняет порядок задач в проекте: перемещает задачу после указанной. Тело запроса:

\begin{lstlisting}[numbers=none, language=json]
{
  "task_id": 7,
  "after_id": 12
}
\end{lstlisting}


Поле after\_id может быть null, тогда задача переносится в начало списка. Ответ -- 204 No Content.

\textbf{Зависимости задач.}

GET /api/projects/\{project\_id\}/dependencies

Запрос не содержит тела. Ответ 200 OK -- массив зависимостей проекта:

\begin{lstlisting}[numbers=none, language=json]
[
  {
    "id": 5,
    "predecessor_id": 7,
    "successor_id": 8,
    "type": "FS"
  }
]
\end{lstlisting}


POST /api/projects/\{project\_id\}/dependencies

Создает зависимость между двумя задачами проекта. Тело запроса:

\begin{lstlisting}[numbers=none, language=json]
{
  "predecessor_id": 7,
  "successor_id": 8,
  "type": "FS"
}
\end{lstlisting}


Ответ 201 Created содержит созданную зависимость. Сервер проверяет принадлежность обеих задач проекту, отсутствие самозависимости, дубля связи, цикла в графе и нарушения дат. Получение отдельной зависимости выполняется через GET /api/projects/\{project\_id\}/dependencies/\{id\}, удаление -- через DELETE по тому же маршруту без тела с ответом 204 No Content.

\textbf{Участники проекта.}

GET /api/projects/\{project\_id\}/members

Запрос не содержит тела. Ответ 200 OK -- массив участников:

\begin{lstlisting}[numbers=none, language=json]
[
  {
    "id": 11,
    "user_id": 3,
    "username": "ivanov.i",
    "full_name": "Иванов Иван Иванович",
    "role": "owner"
  }
]
\end{lstlisting}


POST /api/projects/\{project\_id\}/members

Добавляет пользователя в состав участников проекта. Тело запроса:

\begin{lstlisting}[numbers=none, language=json]
{
  "user_id": 4,
  "role": "editor"
}
\end{lstlisting}


Ответ 201 Created содержит созданную запись участника. Через API нельзя назначить роль owner.

PATCH /api/projects/\{project\_id\}/members/\{id\}

Изменяет роль участника проекта. Тело запроса:

\begin{lstlisting}[numbers=none, language=json]
{
  "role": "viewer"
}
\end{lstlisting}


Ответ 200 OK содержит обновленную запись участника. Удаление участника выполняется через DELETE по тому же маршруту без тела и возвращает 204 No Content; владельца удалить или изменить таким способом нельзя.

DELETE /api/projects/\{project\_id\}/members/leave

Маршрут позволяет текущему пользователю выйти из проекта. Тело запроса не передается. Ответ -- 204 No Content. Владелец проекта выйти не может, вместо этого он должен удалить проект.

\subsection{Проект базы данных}

\subsubsection{ER-модель данных}

Концептуальная модель содержит шесть сущностей \cite{date_db, art_db_methodology}.

Центральная сущность -- проект. Из атрибутов важны название, описание, даты начала и завершения, плюс ссылка на создателя. Внутри проекта находятся задачи. У каждой свое название, описание, даты и порядковый номер для отображения на диаграмме. Между задачами могут быть зависимости одного из четырех типов: финиш-старт, старт-старт, финиш-финиш, старт-финиш. У зависимости два внешних ключа на задачи -- на предшественника и на преемника -- и поле типа связи.

Пользователь системы описывается логином, адресом электронной почты, хешем пароля и ФИО. В одном проекте состоит несколько пользователей в разных ролях -- владелец, редактор, наблюдатель. Эта связь представлена отдельной сущностью «Участие в проекте» с двумя внешними ключами (на пользователя и на проект) и полем роли.

Отдельно от учетных данных у пользователя есть цветовая палитра -- набор цветов и выбранный режим темы. Связь с пользователем один к одному: у каждого аккаунта своя палитра, одна палитра принадлежит одному пользователю.

ER-диаграмма приведена на рисунке~\ref{fig:tp_er}.

\begin{figure}[H]
\centering
\includegraphics[width=\textwidth]{er_model}
\caption{ER-модель данных}
\label{fig:tp_er}
\end{figure}

\subsubsection{Нормализация данных}

Схема приведена к третьей нормальной форме \cite{art_db_normalization}.

Первые две формы выполняются автоматически. Атрибутов-списков и повторяющихся групп в схеме нет -- цвета палитры разложены по отдельным полям, а не упакованы в строку или массив. Все таблицы построены на суррогатном первичном ключе, автоинкрементном поле id, поэтому частичной зависимости от части составного ключа здесь в принципе быть не может. Случаи, где напрашивался составной ключ (пара участник--проект, пара предшественник--преемник), реализованы как ограничения уникальности.

Третья нормальная форма потребовала отдельной проверки. В таблице задач хранятся идентификаторы проекта и автора, а не сами эти объекты. В таблице участников нет ни имени пользователя, ни названия проекта -- только ссылки. Единственная особенность -- производное поле full\_name у пользователя. Оно сделано генерируемым столбцом, который база вычисляет сама по фамилии, имени и отчеству. На третью нормальную форму это не влияет: значение не задается отдельно от других полей строки, оно выводится из них полностью.

\subsubsection{Описание структуры базы данных}

Физическая схема базы состоит из прикладных таблиц системы и служебных таблиц Django. К прикладным относятся users, projects, tasks, project\_members, task\_dependencies и users\_color\_palette. В них хранятся пользователи, проекты, задачи, участие пользователей в проектах, зависимости между задачами и настройки цветовой темы. Остальные таблицы создают встроенные приложения Django: auth\_group, auth\_permission, auth\_group\_permissions, users\_groups, users\_user\_permissions, django\_content\_type, django\_admin\_log, django\_migrations и django\_session.

Прикладные модели используют тип BigAutoField, поэтому первичные ключи id и внешние ключи на них в MySQL имеют тип BIGINT. Текстовые описания проектов и задач хранятся как LONGTEXT, даты -- как DATETIME(6), булевы признаки пользователя -- как TINYINT(1). Связи между прикладными таблицами реализованы внешними ключами. Отношение многие ко многим между пользователями и проектами развернуто в таблицу project\_members с полем роли. Зависимости задач вынесены в task\_dependencies с двумя внешними ключами на tasks -- на предшественника и на преемника. Цветовая палитра связана с пользователем отношением один к одному через поле user\_id.

Схема таблиц с типами полей, первичными и внешними ключами показана на рисунке~\ref{fig:tp_schema}.

\begin{landscape}
\centering
\vspace*{\fill}
\includegraphics[width=\linewidth, height=0.82\textheight, keepaspectratio]{db_struct}

\captionof{figure}{Схема базы данных}
\label{fig:tp_schema}
\vspace*{\fill}
\end{landscape}

Внешние ключи прикладной части настроены либо на каскадное удаление, либо на запрет удаления. При удалении проекта вместе с ним удаляются его задачи, зависимости и записи об участии. При удалении пользователя удаляются его палитра, записи об участии в проектах, связи с группами и индивидуальные разрешения. Пользователя, который числится автором хотя бы одной задачи или создателем хотя бы одного проекта, база удалить не позволит -- у полей creator\_id в таблицах projects и tasks указано on\_delete=RESTRICT.

Структура таблицы users приведена в таблице~\ref{tab:tp_users}.

\begin{xltabular}{\textwidth}{|>{\raggedright\arraybackslash}p{32mm}|C{36mm}|C{36mm}|X|}
\caption{Структура таблицы users}\label{tab:tp_users}\\
\hline
\thead{Поле} & \thead{Тип} & \thead{Обязательность} & \thead{Назначение} \\
\hline
\endfirsthead
\continuecaption{Продолжение таблицы~\thetable}
\thead{Поле} & \thead{Тип} & \thead{Обязательность} & \thead{Назначение} \\
\hline
\finishhead
id & BIGINT & Да & Первичный ключ \\
\hline
password & VARCHAR(128) & Да & Хеш пароля \\
\hline
last\_login & DATETIME(6) & Нет & Дата последнего входа \\
\hline
is\_superuser & TINYINT(1) & Да & Признак суперпользователя \\
\hline
username & VARCHAR(150) & Да & Имя пользователя, уникально \\
\hline
first\_name & VARCHAR(150) & Нет & Имя \\
\hline
last\_name & VARCHAR(150) & Нет & Фамилия \\
\hline
email & VARCHAR(254) & Да & Адрес электронной почты, уникален \\
\hline
is\_staff & TINYINT(1) & Да & Признак доступа к административной панели \\
\hline
is\_active & TINYINT(1) & Да & Признак активного пользователя \\
\hline
date\_joined & DATETIME(6) & Да & Дата регистрации \\
\hline
middle\_name & VARCHAR(150) & Нет & Отчество \\
\hline
full\_name & VARCHAR(452) & Да & Полное имя, вычисляется базой \\
\hline
\end{xltabular}

Структура таблицы projects приведена в таблице~\ref{tab:tp_projects}.

\begin{xltabular}{\textwidth}{|>{\raggedright\arraybackslash}p{32mm}|C{36mm}|C{36mm}|X|}
\caption{Структура таблицы projects}\label{tab:tp_projects}\\
\hline
\thead{Поле} & \thead{Тип} & \thead{Обязательность} & \thead{Назначение} \\
\hline
\endfirsthead
\continuecaption{Продолжение таблицы~\thetable}
\thead{Поле} & \thead{Тип} & \thead{Обязательность} & \thead{Назначение} \\
\hline
\finishhead
id & BIGINT & Да & Первичный ключ \\
\hline
title & VARCHAR(150) & Да & Название проекта \\
\hline
description & LONGTEXT & Нет & Описание \\
\hline
start\_date & DATETIME(6) & Нет & Дата начала проекта \\
\hline
end\_date & DATETIME(6) & Нет & Дата завершения проекта \\
\hline
creator\_id & BIGINT & Да & Внешний ключ на users \\
\hline
\end{xltabular}

Структура таблицы tasks приведена в таблице~\ref{tab:tp_tasks}.

\begin{xltabular}{\textwidth}{|>{\raggedright\arraybackslash}p{32mm}|C{36mm}|C{36mm}|X|}
\caption{Структура таблицы tasks}\label{tab:tp_tasks}\\
\hline
\thead{Поле} & \thead{Тип} & \thead{Обязательность} & \thead{Назначение} \\
\hline
\endfirsthead
\continuecaption{Продолжение таблицы~\thetable}
\thead{Поле} & \thead{Тип} & \thead{Обязательность} & \thead{Назначение} \\
\hline
\finishhead
id & BIGINT & Да & Первичный ключ \\
\hline
title & VARCHAR(150) & Да & Название задачи \\
\hline
description & LONGTEXT & Нет & Описание \\
\hline
start\_date & DATETIME(6) & Нет & Дата начала \\
\hline
end\_date & DATETIME(6) & Нет & Дата завершения \\
\hline
number & INT & Да & Порядковый номер в проекте \\
\hline
creator\_id & BIGINT & Да & Внешний ключ на users \\
\hline
project\_id & BIGINT & Да & Внешний ключ на projects \\
\hline
\end{xltabular}

Пара (project\_id, number) в таблице tasks ограничена уникальностью -- две задачи в одном проекте не могут иметь одинаковый порядковый номер.

Структура таблицы project\_members приведена в таблице~\ref{tab:tp_members}.

\begin{xltabular}{\textwidth}{|>{\raggedright\arraybackslash}p{32mm}|C{36mm}|C{36mm}|X|}
\caption{Структура таблицы project\_members}\label{tab:tp_members}\\
\hline
\thead{Поле} & \thead{Тип} & \thead{Обязательность} & \thead{Назначение} \\
\hline
\endfirsthead
\continuecaption{Продолжение таблицы~\thetable}
\thead{Поле} & \thead{Тип} & \thead{Обязательность} & \thead{Назначение} \\
\hline
\finishhead
id & BIGINT & Да & Первичный ключ \\
\hline
role & VARCHAR(10) & Да & Роль: owner, editor, viewer \\
\hline
project\_id & BIGINT & Да & Внешний ключ на projects \\
\hline
user\_id & BIGINT & Да & Внешний ключ на users \\
\hline
\end{xltabular}

Пара (project\_id, user\_id) ограничена уникальностью -- один пользователь не может состоять в одном проекте дважды.

Структура таблицы task\_dependencies приведена в таблице~\ref{tab:tp_deps}.

\begin{xltabular}{\textwidth}{|>{\raggedright\arraybackslash}p{32mm}|C{36mm}|C{36mm}|X|}
\caption{Структура таблицы task\_dependencies}\label{tab:tp_deps}\\
\hline
\thead{Поле} & \thead{Тип} & \thead{Обязательность} & \thead{Назначение} \\
\hline
\endfirsthead
\continuecaption{Продолжение таблицы~\thetable}
\thead{Поле} & \thead{Тип} & \thead{Обязательность} & \thead{Назначение} \\
\hline
\finishhead
id & BIGINT & Да & Первичный ключ \\
\hline
type & VARCHAR(2) & Да & Тип связи: FS, SS, FF, SF \\
\hline
predecessor\_id & BIGINT & Да & Внешний ключ на tasks (предшественник) \\
\hline
successor\_id & BIGINT & Да & Внешний ключ на tasks (преемник) \\
\hline
\end{xltabular}

Пара (predecessor\_id, successor\_id) ограничена уникальностью. На уровне таблицы запрещено равенство predecessor\_id и successor\_id -- задача не может зависеть от самой себя. Дополнительная проверка на отсутствие циклов выполняется на уровне приложения при создании зависимости.

Структура таблицы users\_color\_palette приведена в таблице~\ref{tab:tp_palette}.

\begin{xltabular}{\textwidth}{|>{\raggedright\arraybackslash}p{32mm}|C{36mm}|C{36mm}|X|}
\caption{Структура таблицы users\_color\_palette}\label{tab:tp_palette}\\
\hline
\thead{Поле} & \thead{Тип} & \thead{Обязательность} & \thead{Назначение} \\
\hline
\endfirsthead
\continuecaption{Продолжение таблицы~\thetable}
\thead{Поле} & \thead{Тип} & \thead{Обязательность} & \thead{Назначение} \\
\hline
\finishhead
id & BIGINT & Да & Первичный ключ \\
\hline
theme\_mode & VARCHAR(16) & Да & Режим темы: system, light, dark, custom \\
\hline
primary\_color & VARCHAR(7) & Да & Базовый цвет палитры \\
\hline
secondary\_\allowbreak color & VARCHAR(7) & Да & Дополнительный цвет палитры \\
\hline
darker\_\allowbreak secondary\_\allowbreak color & VARCHAR(7) & Да & Темный оттенок дополнительного цвета \\
\hline
accent\_\allowbreak color & VARCHAR(7) & Да & Акцентный цвет кнопок и ссылок \\
\hline
task\_\allowbreak color & VARCHAR(7) & Да & Цвет полос задач на диаграмме Ганта \\
\hline
gantt\_\allowbreak grid\_\allowbreak line\_\allowbreak color & VARCHAR(7) & Да & Цвет линий сетки диаграммы \\
\hline
gantt\_\allowbreak current\_\allowbreak time\_\allowbreak color & VARCHAR(7) & Да & Цвет линии текущего времени \\
\hline
error\_color & VARCHAR(7) & Да & Цвет сообщений об ошибках \\
\hline
danger\_color & VARCHAR(7) & Да & Цвет опасных действий \\
\hline
success\_color & VARCHAR(7) & Да & Цвет сообщений об успехе \\
\hline
user\_id & BIGINT & Да & Внешний ключ на users, уникален \\
\hline
app\_\allowbreak background\_\allowbreak color & VARCHAR(7) & Да & Цвет фона приложения \\
\hline
surface\_\allowbreak color & VARCHAR(7) & Да & Цвет поверхностей карточек и панелей \\
\hline
header\_\allowbreak background\_\allowbreak color & VARCHAR(7) & Да & Цвет фона шапки \\
\hline
border\_\allowbreak color & VARCHAR(7) & Да & Цвет границ \\
\hline
primary\_\allowbreak text\_\allowbreak color & VARCHAR(7) & Да & Основной цвет текста \\
\hline
secondary\_\allowbreak text\_\allowbreak color & VARCHAR(7) & Да & Дополнительный цвет текста \\
\hline
inverse\_\allowbreak text\_\allowbreak color & VARCHAR(7) & Да & Цвет текста на акцентных поверхностях \\
\hline
active\_\allowbreak menu\_\allowbreak background\_\allowbreak color & VARCHAR(7) & Да & Цвет фона активного пункта меню \\
\hline
active\_\allowbreak menu\_\allowbreak text\_\allowbreak color & VARCHAR(7) & Да & Цвет текста активного пункта меню \\
\hline
gantt\_\allowbreak arrow\_\allowbreak color & VARCHAR(7) & Да & Цвет стрелок зависимостей \\
\hline
\end{xltabular}

Назначение служебных таблиц Django приведено в таблице~\ref{tab:tp_service_tables}.

\begin{xltabular}{\textwidth}{|>{\raggedright\arraybackslash}p{42mm}|X|}
\caption{Назначение служебных таблиц Django}\label{tab:tp_service_tables}\\
\hline
\thead{Таблица} & \thead{Назначение} \\
\hline
\endfirsthead
\continuecaption{Продолжение таблицы~\thetable}
\thead{Таблица} & \thead{Назначение} \\
\hline
\finishhead
auth\_group & Группы пользователей встроенной системы авторизации Django. \\
\hline
auth\_permission & Разрешения доступа, связанные с типами моделей через content\_type\_id. \\
\hline
auth\_\allowbreak group\_\allowbreak permissions & Связка многие ко многим между группами и разрешениями. \\
\hline
users\_groups & Связка многие ко многим между пользователями и группами. \\
\hline
users\_\allowbreak user\_\allowbreak permissions & Связка многие ко многим между пользователями и индивидуальными разрешениями. \\
\hline
django\_\allowbreak content\_\allowbreak type & Реестр моделей Django, который используется системой разрешений и журналом администратора. \\
\hline
django\_admin\_log & Журнал действий, выполненных через административную панель. \\
\hline
django\_migrations & Список примененных миграций базы данных. \\
\hline
django\_session & Данные пользовательских сеансов; ключ сеанса хранится в cookie браузера. \\
\hline
\end{xltabular}

Для ускорения поиска по часто используемым полям заданы индексы: по title в таблицах проектов и задач, по username и full\_name в таблице пользователей. Первичные ключи и внешние ключи индексируются автоматически.
